\clearpage

\section{总结与展望}

\subsection{论文总结}

本文面向大规模接入物联网场景下的状态更新业务，设计了面向 AoI 的随机接入机制。与传统随机接入主要关注吞吐量或平均碰撞概率不同，本文更加关注平均 AoI 与 AoI 违规概率两个指标。为此，本文将所设计的接入机制引入到时隙 ALOHA 和 802.11 DCF 中，分别构造了等待 ALOHA 协议和等待 CSMA 协议，并对平均 AoI 与 AoI 违规概率进行了联合分析。全文的主要工作和结论可以概括为以下几个方面：
\begin{enumerate}[wide]
    \item \textbf{设计了面向大规模接入场景的等待机制。}针对 AoI 这一指标，本文设计了基于等待的随机接入机制。考虑到该类场景中碰撞可能持续累积并导致有效更新变得稀疏，该机制的基本思路是：节点成功传输后不立即重新参与竞争，而是等待一段时间，并且在碰撞后通过随机退避降低后续的竞争强度。这样处理的目的，不是单纯降低某一次碰撞的概率，而是重新安排节点后续的接入过程，使有效更新在时间轴上的分布更加均衡。

    \item \textbf{等待 ALOHA 的理论分析。}将等待机制引入时隙 ALOHA 框架，构造了等待 ALOHA 协议，建立起了相应的理论模型，并推导出平均 AoI 与 AoI 违规概率的闭式表达式，分析了等待窗口和网络规模对系统性能的影响。

    \item \textbf{等待 CSMA 的理论分析。}将等待机制进一步推广到 802.11 DCF 中，构造了等待 CSMA 协议，并结合载波侦听、退避冻结和变长系统时隙等因素建立起相应的理论模型，推导出平均 AoI 与 AoI 违规概率的闭式表达式，分析等待窗口和网络规模对系统性能的影响。

    \item \textbf{揭示等待机制的作用规律。}仿真结果与理论分析基本一致，说明本文建立的模型能够较准确地刻画系统性能。对比结果表明，在中高负载尤其是海量设备接入的场景下，适度等待并不一定会使信息的时效性发生恶化；相反，通过对节点的接入行为进行适度调节，反而能使有效更新更加均匀地分布在时间轴上，从而改善信息的时效性。
\end{enumerate}

\subsection{研究展望}

尽管本文已经对等待随机接入机制进行了系统的设计与分析，但仍有若干问题值得继续深入。结合本文的模型假设与分析边界，后续工作可以从以下几个方向展开：
\begin{enumerate}[wide]
    \item \textbf{引入更真实的业务到达与排队过程。}本文主要采用按需更新的建模方式，重点分析了接入竞争过程对 AoI 的影响，尚未把随机业务到达、缓存队列和服务过程耦合进统一模型。未来可以在排队论框架下进一步研究等待机制在随机到达流量、有限缓存甚至多业务混合场景中的表现，从而更接近实际网络负载。

    \item \textbf{拓展到更复杂的信道和网络环境。}本文主要关注单跳共享信道环境，对信道衰落、隐藏终端、节点异构以及多跳拓扑等因素考虑有限。后续工作可以把等待机制推广到更复杂的无线环境中，进一步分析不同物理层条件和网络结构对 AoI 表现的影响，并检验现有结论的稳健性。

    \item \textbf{研究参数优化与自适应机制设计。}本文已经得到了平均 AoI 与 AoI 违规概率的解析结果，但对等待窗口、退避参数和系统负载之间的最优关系还缺少更深入的讨论。未来可以在现有闭式基础上建立参数优化问题，并进一步探索依据网络负载或碰撞状态自适应调整等待参数的机制设计。

    \item \textbf{面向实际协议实现开展验证。}本文的结论目前主要建立在理论分析与数值仿真之上。后续若能结合 802.11 实验平台、网络仿真器或软硬件原型系统，对等待机制进行实现和验证，将有助于进一步评估其工程开销、部署难点以及在真实环境中的性能收益。
\end{enumerate}
